\documentclass[11pt]{article}
\usepackage{coursenotes}
\begin{document}
\notesheader{6}{Allocation and policy learning}{From effect estimates to a deployable rule under costs, caps, budgets, and risk}

\begin{decision}
The effect estimates are in. Finance gives a budget, operations a cap on contacts, and management wants a rule simple
enough to explain. Which units get which action, what is an extra yuan of budget worth, how much risk of overspending
is acceptable, and how should a rule be learned and deployed so that it can still be checked?
\end{decision}

\begin{goals}
  \item derive the allocation rule without constraints, under a cap, and under a budget, and interpret the shadow
  price;
  \item decompose net value into the effect's earnings and the sure-thing cost;
  \item price a group treatment-rate constraint by its multiplier;
  \item set a spending reserve and a conservative rule for uncertain effects, and say what each guarantees;
  \item decide under a capacity, where each arm's outcome distribution matters, and set capacity at the critical
  fractile;
  \item place predict-then-optimize, decision-focused learning, and policy learning on one spectrum, and learn a policy
  as weighted classification of doubly robust scores;
  \item deploy a policy so that it can still be evaluated (holdouts, logged probabilities, bandits).
\end{goals}

%=====================================================================================
\sectionone

\subsection{Net value and the unconstrained rule}

Unit $i$ has covariates $X_i$ and decision $d_i \in \{0,1\}$. Treating it earns $m$ per incremental outcome and costs
$c(X_i)$ in expectation, so its \emph{net value} is $v(x) = m\,\tau(x) - c(x)$. Throughout, the value of an assignment is
$\sum_i d_i v(X_i)$ (Section~\ref{sec:sum} says what this requires). The treatment is the one the experiment tested; changing it
(a deeper coupon, no coupon) changes $\tau(x)$, which the test does not measure.

\begin{prop}{Unconstrained rule}{uncon}
Without constraints the optimal assignment is $d_i^\ast = \ind{v(X_i) > 0}$; with a uniform cost $c$ this is
$\ind{\tau(X_i) > c/m}$.
\end{prop}
\begin{proof}
The objective is a sum of terms each depending on one $d_i$; maximize term by term.
\end{proof}

\begin{prop}{Sure-thing decomposition}{surething}
If treatment costs $k$ per unit sent and $a$ per outcome (paid on redemption), so $c(x) = k + a\mu_1(x)$, then
$v(x) = (m - a)\,\tau(x) - a\,\mu_0(x) - k$.
\end{prop}
\begin{proof}
Substitute $\mu_1 = \mu_0 + \tau$ into $m\tau - k - a\mu_1$.
\end{proof}

The first term pays the margin, net of the coupon, on outcomes the treatment causes; the second pays the coupon to units
that would have bought anyway. Ranking by $\tau$ is ranking by $v$ only when costs are uniform. The benefit side needs a
\emph{contrast} ($\tau$, which needs randomization or an identification argument); the cost side needs a
\emph{prediction} ($\mu_0$ or $\mu_1$, fitted on one arm). They have different data requirements and are validated
separately.

\subsection{A cap on the number treated}

\begin{prop}{Caps: rank by net value}{cap}
Under $\sum_i d_i \le B$, an optimal assignment treats the units with the $\min(B, n_+)$ largest values of $v_i$, where
$n_+$ is the number with $v_i > 0$.
\end{prop}
\begin{proof}
Exchange argument: dropping a unit with $v_i \le 0$ never hurts; if a treated $j$ and untreated $i$ have $v_i > v_j$,
swapping them keeps the count and raises the value.
\end{proof}

Without a cap only the sign of $v_i$ matters, and a coarse segment rule that gets signs right is as good as a fine
model. With a binding cap, units that all clear the threshold compete for slots, and order \emph{within} that set earns
money. This is where unit-level estimates add value over segment averages.

\subsection{A budget: the knapsack and its shadow price}

With unit costs $c_i > 0$, benefits $b_i = m\tau_i$, and a budget $C$ on expected cost, solve
$\max_d \sum_i d_i(b_i - c_i)$ subject to $\sum_i d_ic_i \le C$. The \emph{relaxation} allows $d_i \in [0,1]$ (treating a random
share of a segment). Order units by the benefit-cost ratio $r_i = b_i/c_i$; the \emph{greedy} assignment treats units in
that order until the budget runs out or $r_i \le 1$, treating a fraction of the first unit $j$ that does not fit.

\begin{prop}{Greedy is optimal for the relaxation; the shadow price}{knapsack}
Let $\lambda^\ast = r_j - 1$ if the budget binds before the ratio reaches 1, and $0$ otherwise. The greedy assignment solves
the relaxation and equals the rule $d_i = \ind{b_i > (1 + \lambda^\ast)c_i}$ except at the fractional unit. Where the budget
binds, the optimal value $V(C)$ is concave in $C$ with slope $\lambda^\ast$: an extra yuan of budget adds $\lambda^\ast$ of net
value.
\end{prop}
\begin{proof}
For any $\lambda \ge 0$ and feasible $d$,
$\sum_i d_i(b_i - c_i) \le \sum_i d_i\{b_i - (1+\lambda)c_i\} + \lambda C \le \sum_i\{b_i - (1+\lambda)c_i\}^+ + \lambda C$. At
$\lambda^\ast$ the greedy assignment attains both bounds (it treats exactly the positive terms, and spends the budget when
$\lambda^\ast > 0$), so nothing feasible does better. A small budget increase $\delta$ buys $\delta/c_j$ more of unit $j$, worth
$(r_j - 1)\delta = \lambda^\ast\delta$; as $C$ grows the marginal ratio falls, so the slope falls.
\end{proof}

The result has three practical consequences. The budget is spent in ratio order, not in order of $\tau$ or $v$. A
budget can be replaced by a price, since charging each unit's cost at $(1 + \lambda^\ast)$ reproduces the allocation. And
$\lambda^\ast$ is the number to give finance: if the next-best use of the money returns more, the budget belongs there.

The \emph{integer} problem is harder. With costs $6, 5, 5$, net values $12, 9, 9$, and budget $10$, ratio order takes the first (ratio $3.0$) and cannot afford
another, for $12$; the two cheaper units together are worth $18$. Greedy is exact when units are small relative to the
budget or come in blocks (segments); otherwise use an integer solver, and note that greedy is within one unit's value of
the optimum.

\paragraph{Several actions.} With actions $a = 1, \dots, K$ (coupon depths) of benefit $b_i(a)$ and cost $c_i(a)$, the same
Lagrangian argument gives: each unit takes the action maximizing $b_i(a) - (1 + \lambda^\ast)c_i(a)$, or nothing if all are
negative. A deeper coupon is chosen only if its \emph{extra} benefit over the shallower one beats $(1+\lambda^\ast)$ times its
extra cost; tight budgets push toward cheaper actions spread over more units. Contact caps add constraints, each with its
own shadow price.

\subsection{Group constraints and the price of fairness}

A rule that never reads a group label $G$ can still treat groups at very different rates, because effects and costs
differ by group. A \emph{rate constraint} caps the ratio of treatment rates: $\text{rate}_a \le \gamma\,\text{rate}_b$ for
groups $a, b$ of sizes $N_a, N_b$, with $\gamma \ge 1$. The four-fifths guideline of US employment practice (each group's
selection rate at least 80\% of the highest) is $\gamma = 1.25$.

\begin{prop}{Group constraints give group-specific thresholds}{fair}
Under a cap $\sum_i d_i \le B$ and $\sum_{i \in a} d_i/N_a \le \gamma\sum_{i \in b} d_i/N_b$, the relaxation is solved by ranking
units by $v_i$ within each group and treating $i \in a$ if $v_i > \lambda^\ast + \eta^\ast/N_a$ and $i \in b$ if
$v_i > \lambda^\ast - \gamma\eta^\ast/N_b$ (other groups at $\lambda^\ast$), where $\lambda^\ast, \eta^\ast \ge 0$ are the multipliers of
the cap and the rate constraint. Where both bind, the optimal value $V(\gamma)$ rises with $\gamma$ at rate
$\eta^\ast\,\text{rate}_b$.
\end{prop}
\begin{proof}
The Lagrangian $\sum_i d_iv_i - \lambda(\sum_i d_i - B) - \eta(\sum_{i \in a} d_i/N_a - \gamma\sum_{i \in b} d_i/N_b)$ is a sum of
terms in single $d_i$, each maximized by the stated threshold. The weak-duality argument of Result~\ref{prop:knapsack} shows
that an assignment meeting both constraints with complementary slackness is optimal, and the envelope theorem gives
$dV/d\gamma = \partial/\partial\gamma$ of the Lagrangian, $\eta^\ast\sum_{i \in b} d_i/N_b$.
\end{proof}

In effect, a constraint on rates sets different thresholds for the two groups. The group the unconstrained rule
favors faces a higher bar, the other a lower one, and $\eta^\ast$ sets the gap. The \emph{price of fairness}
$V(\infty) - V(\gamma)$ is the value given up, and it is the value of the units swapped.

Take 10{,}000 customers in two groups of 5{,}000 and a cap of 1{,}000. Group $a$ has 750 customers at $v = $ ¥6 and the rest
at ¥1; group $b$ has 250 at ¥6, 1{,}000 at ¥3, and the rest at ¥0.5. The unconstrained list is the 1{,}000 customers at ¥6,
with rates of 15\% and 5\% (three to one) and a value of ¥6{,}000. With $\gamma = 1.25$ the rates satisfy
$5{,}000(\text{rate}_b + 1.25\,\text{rate}_b) = 1{,}000$, so group $b$ gets $444.4$ and group $a$ $555.6$ (in whole customers,
445 and 555). The list swaps about 195 customers at ¥6 in $a$ for 195 at ¥3 in $b$. Its value falls to about ¥5{,}415, a
price of about ¥585. The thresholds are ¥6 in $a$ and ¥3 in $b$; solving the two threshold equations gives $\eta^\ast \approx$ ¥6{,}667 and $\lambda^\ast \approx$ ¥4.67, so tightening $\gamma$ from 1.25 to 1.20 costs about
$0.05 \times 6{,}667 \times 0.089 \approx$ ¥30.

Which constraint to impose is a value judgment rather than an estimation problem. Parity of treatment rates, of expected
benefit, and of outcomes are different constraints with different prices, and they can conflict. Writing the group label
into the rule, as group-specific thresholds do, may itself be restricted where the label is legally protected. The
multiplier tells management what a constraint costs; whether to pay it is their decision.

\subsection{Spending risk}

The budget above constrains \emph{expected} cost. With redemption $R_i \sim \text{Bernoulli}(s_i)$, $s_i = \mu_1(X_i)$, unit $i$
costs $k + aR_i$.

\begin{prop}{Mean, variance, and the reserve}{reserve}
If redemptions are independent, realized cost $K = \sum_i d_i(k + aR_i)$ has $\E[K] = \sum_i d_ic_i$ and
$\Var(K) = a^2\sum_i d_is_i(1-s_i)$. Under the normal approximation, $\Pr(K > C) \le \alpha$ if and only if
$\E[K] + z_\alpha\sqrt{\Var(K)} \le C$.
\end{prop}
\begin{proof}
Linearity gives the mean; independence makes variances add, each $a^2s_i(1-s_i)$. With $K$ approximately normal,
$\Pr(K > C) = 1 - \Phi\{(C - \E K)/\text{sd}(K)\}$.
\end{proof}

The \emph{reserve} $z_\alpha\,\text{sd}(K)$ grows like the square root of the list while expected cost grows linearly, so it is a
small share of a large campaign. It covers Bernoulli luck, not a wrong cost model or a common shock (a heatwave) that moves
every redemption together.

\subsection{Uncertain effects: the conservative rule}

If $\hat v(x)$ is approximately $N(v(x), \SE(x)^2)$, the \emph{conservative rule} treats when
$\hat v(x) - z_\alpha\SE(x) > 0$.

\begin{prop}{What the conservative rule guarantees}{conservative}
For any $x$ with $v(x) \le 0$, the conservative rule treats it with probability at most $\alpha$.
\end{prop}
\begin{proof}
If $v \le 0$, $\Pr\{\hat v - z_\alpha\SE > 0\} \le \Pr\{\hat v - z_\alpha\SE > v\} = \alpha$.
\end{proof}

The guarantee is one-sided. It limits paying for losers and says nothing about profitable cells left off because
their estimates are noisy. $z_\alpha$ expresses a risk preference, not something the data supply. Bound the net value, not
$\tau$, because the cost is estimated too. With arms estimated independently, $\widehat\Var(\hat v) = (m - a)^2\widehat\Var(\hat\mu_1) + m^2\widehat\Var(\hat\mu_0)$ (Meeting~2). Optimizing
on point estimates also suffers the \emph{optimizer's curse}. Units chosen because their estimates are high are
disproportionately those whose estimates are too high, so the realized value of an optimized list falls short of its
estimate. Evaluate the frozen list on held-out data (Meeting~5).

\subsection{When the mean is not enough: capacity}

Net value $m\tau(x) - c(x)$ is linear in the outcome. When the payoff is not (a kitchen that can cook only so many meals,
a store with a fixed number of staff), the average effect no longer decides. Make the unit the holder of the capacity: a
partner kitchen with daily orders $Y$, capacity $q$, margin $p$ per order served, and cost $g$ per order turned away
(refunds, goodwill). A day is worth $\pi(y) = p\min(y, q) - g(y - q)^+$, and running the promotion costs $c$.

\begin{prop}{A nonlinear payoff needs each arm's distribution}{nonlin}
The promotion's net value at $x$ is $v(x) = \E[\pi(Y(1)) \mid X = x] - \E[\pi(Y(0)) \mid X = x] - c$. In a randomized experiment
each term is identified from one arm, $\E[\pi(Y(d)) \mid X = x] = \E[\pi(Y) \mid D = d, X = x]$. So $v(x)$ depends on the
conditional distributions of $Y(1)$ and $Y(0)$ separately; it does not depend on their joint distribution, and it is not
a function of $\tau(x)$ alone unless $\pi$ is linear ($q = \infty$ gives $v = p\tau(x) - c$).
\end{prop}
\begin{proof}
Linearity of expectation splits $\E[\pi(Y(1)) - \pi(Y(0)) \mid X]$ into one term per arm; independence of $D$ and $Y(d)$ given $X$
and consistency give the identification, as in Meeting~3.
\end{proof}

Two kitchens with baseline 200 orders, capacity 210, $p = $ ¥25, and $g = $ ¥15. Kitchen A's promoted demand is 215 for
certain; kitchen B's is 230 or 200 with probability one half each. Both have $\tau = 15$. For A,
$v + c = 25 \times 10 - 15 \times 5 = 175$. For B, $v + c = \tfrac12(25 \times 10 - 15 \times 20) + \tfrac12 \times 0 = -25$. The two
kitchens have the same CATE but call for opposite decisions. What the decision does \emph{not} need is the distribution of individual effects $Y(1) - Y(0)$, which no
experiment identifies, since no kitchen shows both outcomes on the same day.

\paragraph{Choosing the capacity.} Suppose the kitchen prepares $q$ meals before demand is known, losing $c_u$ per order it
cannot fill and $c_o$ per meal prepared and wasted. Let $F_1(\cdot \mid x)$ be the distribution function of $Y(1)$ given $X = x$.

\begin{prop}{The critical fractile}{fractile}
The expected loss $c_u\E[(Y(1) - q)^+ \mid x] + c_o\E[(q - Y(1))^+ \mid x]$ is minimized at the $\kappa$-quantile of $F_1(\cdot \mid x)$,
$\kappa = c_u/(c_u + c_o)$.
\end{prop}
\begin{proof}
The derivative in $q$ is $-c_u\{1 - F_1(q \mid x)\} + c_oF_1(q \mid x)$: increasing, and zero at $F_1(q \mid x) = \kappa$.
\end{proof}

With $c_u = $ ¥40 (margin plus goodwill) and $c_o = $ ¥10, $\kappa = 0.8$. Kitchen B should prepare 230, not the 215 that
``baseline plus CATE'' suggests; the expected loss falls from ¥375 to ¥150 a day. Three further points. (i) $\kappa$ comes from the
business, not from a risk preference: it prices the randomness of the outcome. The $z_\alpha$ of the conservative rule
guards against error in the \emph{estimate} and remains a preference; a full treatment uses both. (ii) The quantile of
$Y(1)$ given $x$ is a prediction on the treated arm, fitted by minimizing the \emph{pinball loss}
$\rho_\kappa(u) = u\{\kappa - \ind{u < 0}\}$ (quantile regression, or the quantile forest of Meeting~4); like the cost side of
Result~\ref{prop:surething}, it needs one arm. (iii) A \emph{conditional quantile treatment effect},
$F_1^{-1}(\kappa \mid x) - F_0^{-1}(\kappa \mid x)$, compares the two arms' quantiles. It is not the quantile of individual effects,
and neither decision above uses it.

\subsection{From predict-then-optimize to policy learning}

Every rule so far is \emph{predict-then-optimize}. It estimates $\tau$ (and $c$), then hands the estimates to an
optimizer (a threshold, a sort, a knapsack). The estimates are trained for accuracy; the decision is judged by value. Three routes
differ in how much of the training they hand to the decision.
\begin{enumerate}
  \item \emph{Predict, then optimize} (plug-in). Fit $\hatt$ by an estimation loss and apply the optimizer. Modular and
  reusable, but accuracy is bought everywhere, including where the decision is never in doubt.
  \item \emph{Decision-focused learning.} Keep the optimizer and train the estimate on the value of the decision it
  induces, choosing $\theta$ to maximize $\frac1n\sum_i \pi_\theta(X_i)\Gamma_i$, where $\pi_\theta$ is the optimizer applied to
  $\hatt_\theta$ and $\Gamma_i$ is the score of Result~\ref{prop:ewm} below. The objective jumps as the optimizer's output
  changes, so practice uses smooth surrogates (the SPO+ loss of Elmachtoub and Grigas (2022) for linear objectives) or
  randomized optimizers. It is worth the extra effort when the optimizer couples units (a budget, a capacity, a menu), so that
  a unit's decision is not a simple function of its own estimate.
  \item \emph{Policy learning.} Drop the estimate and search a class $\Pi$ of rules, such as depth-two trees, for the one with
  the highest estimated value. A restricted class is easy to explain, audit, and deploy.
\end{enumerate}
Route 2 is route 3 with the class $\Pi = \{\text{optimizer} \circ \hatt_\theta\}$. One difference from textbook
predict-then-optimize: there the training labels (realized demands, costs) are observed noisy draws of their conditional
mean. Here the label $\tau_i$ is never observed; the doubly robust score takes its place, as a noisy draw whose
conditional mean is the net value. With that substitution, predict-then-optimize methods carry over.

\begin{prop}{Empirical welfare maximization with doubly robust scores}{ewm}
Let $\Gamma_i = m\,\phi_i - c(X_i)$, where $\phi_i$ is the cross-fitted doubly robust score of Meeting~3. Then
$\E[\Gamma_i \mid X_i] = v(X_i)$, so $\frac1n\sum_i \pi(X_i)\Gamma_i$ is unbiased for the gain $G(\pi)$ of any fixed policy, and
\[
  \hat\pi = \arg\max_{\pi \in \Pi} \frac1n\sum_{i=1}^n \pi(X_i)\,\Gamma_i .
\]
\end{prop}
\begin{proof}
$\E[\phi_i \mid X_i] = \tau(X_i)$ when either the outcome models or the propensity are correct (Meeting~3), so
$\E[\Gamma_i \mid X_i] = m\tau(X_i) - c(X_i)$; then apply the policy-value identity of Meeting~5.
\end{proof}

\begin{prop}{Policy learning is weighted classification}{wclass}
For any policy $\pi$, with $\Gamma_i^- = \max(-\Gamma_i, 0)$,
\[
  \sum_{i=1}^n \pi(X_i)\Gamma_i = \sum_{i=1}^n |\Gamma_i|\,\ind{\pi(X_i) = \ind{\Gamma_i > 0}} - \sum_{i=1}^n \Gamma_i^- .
\]
Maximizing estimated value is minimizing the misclassification of the sign of $\Gamma_i$, each unit weighted by $|\Gamma_i|$.
\end{prop}
\begin{proof}
If $\Gamma_i > 0$, $\pi\Gamma_i = |\Gamma_i|\ind{\pi = 1}$. If $\Gamma_i \le 0$, $\pi\Gamma_i = -|\Gamma_i|\pi = |\Gamma_i|\ind{\pi = 0} - |\Gamma_i|$.
The last sum does not depend on $\pi$.
\end{proof}

Any classifier that accepts sample weights becomes a policy learner (\emph{outcome-weighted learning}). This makes the
contrast with route 1 precise: an estimation loss penalizes errors on every unit, while the decision loss penalizes only errors
that flip a decision, weighted by how much is at stake. Later in the course, a population version of this result shows that regret counts
only misclassified units. Under a budget, replace $\Gamma_i$ by $\Gamma_i - \lambda^\ast c(X_i)$ (Result~\ref{prop:knapsack}).

A \emph{policy tree} is the case where $\Pi$ is the set of trees of depth one or two, searched exhaustively. The \emph{regret}
of $\hat\pi$, the value it gives up relative to the best policy in $\Pi$, shrinks at a rate of order $\sqrt{\text{complexity}(\Pi)/n}$
with doubly robust scores (Athey and Wager 2021). A richer class can come closer to the ideal rule but needs more data
to find its best member, and depth-two trees are a common compromise. The selection warning of Meeting~5 applies here
too. The learned policy's in-sample value is optimistic, so evaluate it on fresh data.

\subsection{Deploying so that it can still be checked}

A deployed policy is a hypothesis about where value lies. Keep a small random \emph{holdout} of eligible units under the
old policy or untreated. It keeps measuring value, detects drift, and generates experimental data for the next round. A
policy deployed without one destroys the data needed to evaluate it. When the rule itself randomizes (for instance, a
bandit), log each unit's assignment probability: the logs are then experimental data with known, varying propensities,
and the IPW identity of Meeting~5 applies with those probabilities.

\paragraph{Bandits.} A fixed test spends half its traffic on the worse arm until the end. A \emph{multi-armed bandit} shifts
traffic toward arms that look better as data arrive: \emph{Thompson sampling} assigns each new unit to an arm with
probability equal to the current posterior probability that the arm is best. Bandits suit many short-lived decisions
(creative variants); one-off strategic decisions are better served by a clean test. Sequential rules, where today's
action depends on yesterday's response, need more care. The response is both a result of the first action and a reason
for the next, so adjusting for it by ordinary regression biases the effect of the first action.

\subsection{What the sum assumes}\label{sec:sum}

Writing the value of a list as $\sum_i d_iv(X_i)$ assumes no interference and one well-defined treatment. It fails when
treated units tell untreated ones, draw demand from them, or share a capacity. Then the objective is not
separable and none of the sorting results holds. Check also that the list does not rank cells the experiment never
contained (no support): the optimizer cannot tell a measured number from an extrapolated one.

\begin{pitfall}[allocation errors]
Ranking by $\tau$ when costs differ by unit; spending the budget in order of effect rather than ratio; planning capacity
at baseline plus the CATE; treating a fairness constraint as free; reporting the optimizer's own estimate of the list's
value; deploying a rule with no holdout; ignoring spillovers when a list is scaled up.
\end{pitfall}

%=====================================================================================
\sectiontwo

\paragraph{Setting.} Pujiang Delivery can send a ¥5 or a ¥10 coupon. Six segments (A--F) of 10{,}000 customers were randomized
across the three arms last month. Per customer offered, the experiment gives incremental margin $b$ and expected coupon
cost $c$ (face value times redemption rate, including redemptions by customers who would have ordered anyway), in yuan.
\begin{center}
\begin{tabular}{lcccccc}
\toprule
& \multicolumn{3}{c}{¥5 coupon} & \multicolumn{3}{c}{¥10 coupon} \\
\cmidrule(lr){2-4}\cmidrule(lr){5-7}
Segment & $b$ & $c$ & $b - c$ & $b$ & $c$ & $b - c$ \\
\midrule
A & 6.0 & 2.0 & 4.0 & 9.0 & 5.0 & 4.0 \\
B & 4.0 & 2.0 & 2.0 & 7.0 & 4.5 & 2.5 \\
C & 3.0 & 2.5 & 0.5 & 4.0 & 5.0 & $-1.0$ \\
D & 1.5 & 2.5 & $-1.0$ & 2.0 & 5.0 & $-3.0$ \\
E & 2.0 & 1.0 & 1.0 & 2.5 & 2.5 & 0.0 \\
F & 0.5 & 1.5 & $-1.0$ & 0.8 & 3.0 & $-2.2$ \\
\bottomrule
\end{tabular}
\end{center}

\paragraph{Step 1: no budget.} Each segment takes its best positive action: A ¥5 (same net value as ¥10 at lower cost), B ¥10,
C ¥5, E ¥5. Spend $10{,}000 \times (2.0 + 4.5 + 2.5 + 1.0) = $ ¥100{,}000; profit ¥80{,}000.

\paragraph{Step 2: a ¥50{,}000 budget.} Net value per yuan of cost: A at ¥5, $4.0/2.0 = 2.0$; B and E at ¥5, $1.0$; C at ¥5, $0.2$;
upgrading B from ¥5 to ¥10, $(2.5 - 2.0)/(4.5 - 2.0) = 0.2$. Filling from the top, A, B, and E at ¥5 use exactly ¥50{,}000 and earn
¥70{,}000; a linear-programming solver returns the same. The next uses of money return ¥0.20 per yuan, so
$\lambda^\ast = 0.2$: an extra ¥10{,}000 would add about ¥2{,}000. Halving the budget costs only ¥10{,}000 of profit because the
dropped spending was the least productive. (At exactly ¥50{,}000, \emph{removing} a yuan would cost ¥1.00, since B or E would be
cut; the shadow price differs on the two sides of a kink.)

\paragraph{Step 3: the price rule.} With $\lambda = 0.2$, B scores $2.0 - 0.2 \times 2.0 = 1.6$ at ¥5 and $2.5 - 0.2 \times 4.5 = 1.6$ at ¥10:
indifferent, which is why $0.2$ is where the budget starts to bind on B's upgrade. C's ¥5 coupon scores
$0.5 - 0.2 \times 2.5 = 0$.

\paragraph{Step 4: noise.} C's net value of ¥0.5 per customer has SE ¥0.4. A 90\% lower bound is $0.5 - 1.28 \times 0.4 = -0.01$, so
the conservative rule skips C even with unlimited budget, at little expected cost since C is marginal.

\paragraph{Step 5: a policy tree within B.} Eight customers from the experiment with cross-fitted DR scores $\Gamma_i$ for the ¥5
coupon (net of cost):
\begin{center}
\begin{tabular}{ccccccccc}
\toprule
Customer & 1 & 2 & 3 & 4 & 5 & 6 & 7 & 8 \\
\midrule
New customer? & yes & yes & yes & yes & no & no & no & no \\
High spend? & yes & yes & no & no & yes & yes & no & no \\
$\Gamma_i$ & 3.0 & 1.0 & 2.0 & $-1.0$ & $-2.0$ & 0.5 & $-0.5$ & $-1.5$ \\
\bottomrule
\end{tabular}
\end{center}
The estimated gain per customer, $\tfrac18\sum_i\pi(X_i)\Gamma_i$, is $1.5/8 = 0.19$ for treating everyone, $5.0/8 = 0.63$
for new customers, and $2.5/8 = 0.31$ for high spenders. The best depth-one tree treats new customers; depth two cannot improve (cell sums $4.0, 1.0, -1.5,
-2.0$). By Result~\ref{prop:wclass} the same search is a weighted classification, which labels each customer by the sign
of $\Gamma_i$ and weights it by $|\Gamma_i|$. ``New customers'' misclassifies customers 4 and 6 (weight $1.5$); ``everyone''
misclassifies 4, 5, 7, and 8 (weight $5.0$). The difference, $3.5$, is the difference in summed value, $5.0 - 1.5$. Eight
customers are far too few to trust; the example only illustrates the arithmetic.

\paragraph{Step 6: deployment.} Pujiang Delivery deploys with 5\% of each segment held out at random under no coupon, so next month's
analysis can re-estimate $b$ and $c$.

\begin{exercise}[computation]
Finance offers to raise the budget from ¥50{,}000 to ¥75{,}000. (a) What is the new allocation and profit? (b) Is the extra ¥25{,}000
worth it if marketing's alternative use of money returns ¥0.30 per yuan?
\end{exercise}
\answer{(a) After A, B, E at ¥5 (¥50{,}000), the next options both return $0.2$: C at ¥5 costs ¥25{,}000 for ¥5{,}000 net; upgrading B costs
¥25{,}000 for ¥5{,}000. Either (or a mix) gives profit ¥75{,}000. (b) The extra money returns $5{,}000/25{,}000 = 0.2$ per yuan, below
marketing's $0.30$: give it to marketing.}

\begin{exercise}[judgment]
The growth team wants to send the ¥10 coupon to segment A ``because its incremental margin, ¥9.0, is the highest of any option''.
What is wrong with the argument, and what would change your mind?
\end{exercise}
\answer{Incremental margin is not net value: A's ¥10 coupon costs ¥5.0 and nets ¥4.0, the same as the ¥5 coupon at less than half the
cost, so under any binding budget the ¥5 coupon is better (ratio 2.0 vs 0.8). Without a budget they tie. Evidence that A's ¥10
coupon has a larger net value (a lower redemption subsidy, or a longer-lasting effect) from a new test would change the answer.}

\begin{exercise}[computation]
A third partner kitchen has baseline 200 orders, capacity 210, $p = $ ¥25, $g = $ ¥15, and a promotion costing ¥60 a day. A test
shows promoted demand of 205, 215, or 235 orders, each with probability one third. (a) Compute $\tau$ and the linear rule's
verdict $p\tau - c$. (b) Compute the correct net value (Result~\ref{prop:nonlin}). (c) If instead the kitchen can prepare meals in
advance with $c_u = $ ¥40 and $c_o = $ ¥10, how many should it prepare for a promoted day?
\end{exercise}
\answer{(a) $\tau = 218.3 - 200 = 18.3$; $25 \times 18.3 - 60 = 398$: promote. (b) The incremental day values are $25 \times 5 = 125$,
$25 \times 10 - 15 \times 5 = 175$, and $25 \times 10 - 15 \times 25 = -125$; their mean is $58.3$, so $v = 58.3 - 60 = -1.7$: do not promote.
The tail day above capacity costs more than the average lift earns. (c) $\kappa = 40/50 = 0.8$; $F_1(215) = 2/3 < 0.8 \le F_1(235) = 1$,
so prepare 235.}

\begin{exercise}[judgment]
Pujiang Delivery's effect estimates are noisier for group $b$, which was smaller in the experiment. The team applies the conservative
rule and the rate constraint of Result~\ref{prop:fair} together. What happens to the rates, to $\eta^\ast$, and what would you change
for next month?
\end{exercise}
\answer{Wider SEs push more group-$b$ customers below the conservative bound, so the unconstrained list treats group $b$ even less;
the rate constraint binds harder, more swaps are forced, and $\eta^\ast$ (the price of fairness) rises. Part of that price is a price of
ignorance, not of fairness. Fix it at the source: stratify next month's experiment to over-sample group $b$, which narrows its SEs
at little cost.}

%=====================================================================================
\sectionthree

\begin{extension}[Decision-focused learning in practice]
The value of a decision is piecewise constant in the estimates, so its gradient is zero almost everywhere. Three fixes:
convex surrogates such as SPO+ (Elmachtoub and Grigas 2022), which bound the regret of a linear optimizer; smoothing the
optimizer by adding noise and averaging its output; and, for threshold and sort rules, the weighted classification of
Result~\ref{prop:wclass}. Mandi et al.\ (2024) survey the methods and benchmarks. With causal labels, every surrogate is fed
doubly robust scores in place of realized outcomes, and the cross-fitting rules of Meeting~3 apply.
\end{extension}

\begin{extension}[Distributional effects]
Quantile treatment effects compare the arms' outcome distributions at a fixed quantile; they are identified in an experiment
and estimable by weighting (Firpo 2007) or by quantile forests. Ban and Rudin (2019) develop data-driven newsvendor rules
that fit the critical quantile directly from features, the one-arm version of Result~\ref{prop:fractile}.
\end{extension}

\begin{extension}[Theory of policy learning]
Kitagawa and Tetenov (2018) bound the regret of empirical welfare maximization with known propensities; Athey and Wager (2021)
extend it to doubly robust scores with estimated nuisances. The \texttt{policytree} package implements exact tree search.
\end{extension}

\subsection*{Annotated reading}
\begin{readinglist}
\reading{Kitagawa, Toru, and Aleksey Tetenov. 2018. ``Who Should Be Treated? Empirical Welfare Maximization Methods
for Treatment Choice.'' \emph{Econometrica} 86 (2): 591--616. \url{https://doi.org/10.3982/ECTA13288}}{The case for
choosing a rule from a restricted class, with regret bounds}{Sections 1--3}{3}
\reading{Athey, Susan, and Stefan Wager. 2021. ``Policy Learning with Observational Data.'' \emph{Econometrica} 89
(1): 133--161. \url{https://doi.org/10.3982/ECTA15732}}{Doubly robust scores and policy trees}{Sections 1--3}{3}
\reading{Elmachtoub, Adam N., and Paul Grigas. 2022. ``Smart `Predict, then Optimize'.'' \emph{Management Science} 68
(1): 9--26. \url{https://doi.org/10.1287/mnsc.2020.3922}}{Training predictions on the decisions they induce; the SPO+
loss}{Sections 1--4}{3}
\reading{Mandi, Jayanta, James Kotary, Senne Berden, Maxime Mulamba, Victor Bucarey, Tias Guns, and Ferdinando
Fioretto. 2024. ``Decision-Focused Learning: Foundations, State of the Art, Benchmark and Future Opportunities.''
\emph{Journal of Artificial Intelligence Research} 80: 1623--1701. \url{https://doi.org/10.1613/jair.1.15320}}{A
survey of decision-focused methods, with a benchmark of 11 of them}{Introduction}{2}
\reading{Zhao, Yingqi, Donglin Zeng, A. John Rush, and Michael R. Kosorok. 2012. ``Estimating Individualized Treatment
Rules Using Outcome Weighted Learning.'' \emph{Journal of the American Statistical Association} 107 (499):
1106--1118. \url{https://doi.org/10.1080/01621459.2012.695674}}{Policy learning as weighted
classification}{Sections 1--3}{2}
\reading{Firpo, Sergio. 2007. ``Efficient Semiparametric Estimation of Quantile Treatment Effects.''
\emph{Econometrica} 75 (1): 259--276. \url{https://doi.org/10.1111/j.1468-0262.2007.00738.x}}{Quantile treatment
effects estimated by weighting}{Introduction}{3}
\reading{Ban, Gah-Yi, and Cynthia Rudin. 2019. ``The Big Data Newsvendor: Practical Insights from Machine Learning.''
\emph{Operations Research} 67 (1): 90--108. \url{https://doi.org/10.1287/opre.2018.1757}}{Fitting the critical
quantile directly from features}{Introduction}{2}
\reading{Kasy, Maximilian, and Rediet Abebe. 2021. ``Fairness, Equality, and Power in Algorithmic Decision-Making.''
In \emph{Proceedings of the 2021 ACM Conference on Fairness, Accountability, and Transparency (FAccT '21)}, 576--586.
\url{https://doi.org/10.1145/3442188.3445919}}{What fairness constraints do and do not protect, from an economist's
view}{All}{1}
\reading{Russo, Daniel J., Benjamin Van Roy, Abbas Kazerouni, Ian Osband, and Zheng Wen. 2018. ``A Tutorial on
Thompson Sampling.'' \emph{Foundations and Trends in Machine Learning} 11 (1): 1--96.
\url{https://doi.org/10.1561/2200000070}}{Bandits from first principles}{Ch.\ 1--4}{2}
\reading{Handout H3, \emph{Sequential decisions}}{Dynamic treatment regimes, Q-learning, G-estimation, and bandit
logs}{All}{2}
\end{readinglist}

\end{document}
